\chapter{Lesson 12: Short Forms and Informal Speech}

\section{Long forms and short forms}

So far every sentence has ended in a \textbf{long form} (polite form): \ja{です},
\ja{ます}. Each long form has a corresponding \textbf{short form} (plain form). Short
forms are used:
\begin{enumerate}
  \item as the casual speech style with friends and family;
  \item inside longer sentences (before \ja{と思います}, \ja{って}, \ja{から}, noun
        modifiers, \ldots), \emph{even in polite speech}.
\end{enumerate}
The dictionary form we have been using is the present affirmative short form.

\begin{center}
\begin{tikzpicture}[x=1cm,y=1cm, lv/.style={draw, rounded corners, minimum width=5.8cm, minimum height=0.8cm, font=\small}]
  \draw[-Stealth, thick] (-3.4,-0.3) -- (-3.4,3.3) node[above, font=\small] {formality};
  \node[lv, fill=purple!10] at (0,3) {keigo (honorific/humble)};
  \node[lv, fill=blue!10] at (0,2) {polite: \ja{たべます}, \ja{がくせいです}};
  \node[lv, fill=green!10] at (0,1) {short/plain: \ja{たべる}, \ja{がくせいだ}};
  \node[lv, fill=orange!10] at (0,0) {casual speech: \ja{たべる？}, \ja{がくせい？}};
  \node[font=\scriptsize, anchor=west] at (3.1,3) {customers, superiors (later courses)};
  \node[font=\scriptsize, anchor=west] at (3.1,2) {teachers, strangers, class};
  \node[font=\scriptsize, anchor=west] at (3.1,1) {writing, inside sentences};
  \node[font=\scriptsize, anchor=west] at (3.1,0) {friends, family, juniors};
\end{tikzpicture}
\end{center}

\section{Short forms of verbs: present}

\subsection{Present affirmative}
The dictionary form: \ja{たべる}, \ja{のむ}, \ja{する}, \ja{くる}.

\subsection{Present negative: \ja{〜ない}}
\begin{itemize}
  \item \ja{る}-verbs: drop \ja{る}, add \ja{ない}: \ja{たべる} $\to$ \ja{たべない}.
  \item \ja{う}-verbs: change final \emph{u} to \emph{a}, add \ja{ない}:
        \ja{のむ} $\to$ \ja{のまない}, \ja{かく} $\to$ \ja{かかない}.
        Verbs ending in \ja{う} use \ja{わ}: \ja{かう} $\to$ \ja{かわない}.
  \item Irregular: \ja{する} $\to$ \ja{しない}, \ja{くる} $\to$ \ja{こない}.
  \item Special: \ja{ある} $\to$ \ja{ない} (not *\ja{あらない}).
\end{itemize}

\section{Short forms of verbs: past}

\subsection{Past affirmative: \ja{〜た}}
Take the \ja{て}-form and change \ja{て}$\to$\ja{た}, \ja{で}$\to$\ja{だ}. All the
\ja{て}-form rules (Lesson 11) carry over.
\begin{center}
\begin{tabular}{lll}
\toprule
Dictionary & \ja{て}-form & Past short \\
\midrule
\ja{たべる} & \ja{たべて} & \ja{たべた} \\
\ja{かう} & \ja{かって} & \ja{かった} \\
\ja{よむ} & \ja{よんで} & \ja{よんだ} \\
\ja{かく} & \ja{かいて} & \ja{かいた} \\
\ja{およぐ} & \ja{およいで} & \ja{およいだ} \\
\ja{はなす} & \ja{はなして} & \ja{はなした} \\
\ja{いく} & \ja{いって} & \ja{いった} \\
\ja{する} & \ja{して} & \ja{した} \\
\ja{くる} & \ja{きて} & \ja{きた} \\
\bottomrule
\end{tabular}
\end{center}

\subsection{Past negative: \ja{〜なかった}}
Take the \ja{ない}-form and replace \ja{ない} with \ja{なかった}:
\ja{たべなかった}, \ja{のまなかった}, \ja{しなかった}, \ja{こなかった}, \ja{なかった}.

\section{Summary table: long vs.\ short}

\begin{center}
\begin{tabular}{lllll}
\toprule
 & present aff. & present neg. & past aff. & past neg. \\
\midrule
verb (long) & \ja{よみます} & \ja{よみません} & \ja{よみました} & \ja{よみませんでした} \\
verb (short) & \ja{よむ} & \ja{よまない} & \ja{よんだ} & \ja{よまなかった} \\
\midrule
\ja{い}-adj (long) & \ja{さむいです} & \ja{さむくないです} & \ja{さむかったです} & \ja{さむくなかったです} \\
\ja{い}-adj (short) & \ja{さむい} & \ja{さむくない} & \ja{さむかった} & \ja{さむくなかった} \\
\midrule
\ja{な}-adj (long) & \ja{ひまです} & \ja{ひまじゃないです} & \ja{ひまでした} & \ja{ひまじゃなかったです} \\
\ja{な}-adj (short) & \ja{ひまだ} & \ja{ひまじゃない} & \ja{ひまだった} & \ja{ひまじゃなかった} \\
\midrule
noun (long) & \ja{がくせいです} & \ja{がくせいじゃないです} & \ja{がくせいでした} & \ja{がくせいじゃなかったです} \\
noun (short) & \ja{がくせいだ} & \ja{がくせいじゃない} & \ja{がくせいだった} & \ja{がくせいじゃなかった} \\
\bottomrule
\end{tabular}
\end{center}

\begin{important}[\ja{い}-adjectives: just drop \ja{です}]
For \ja{い}-adjectives the short form is the long form minus \ja{です}. Never add
\ja{だ} to an \ja{い}-adjective: *\ja{さむいだ} is wrong.
\end{important}

\section{Casual speech}

In conversation among friends:
\begin{itemize}
  \item Use short forms at the end of sentences.
  \item Questions drop \ja{か} and use rising intonation: \ja{あしたひま？} ``Free
        tomorrow?'' (\ja{か} in casual speech sounds rough.)
  \item \ja{だ} is usually dropped at the end of a statement or question, especially by
        women; \ja{だよ}, \ja{だね} are common: \ja{いいてんきだね}.
  \item Particles \ja{は}, \ja{を}, \ja{が} are often dropped: \ja{コーヒー、のむ？}
  \item \ja{うん} ``yeah'' and \ja{ううん} ``nah'' replace \ja{はい}/\ja{いいえ}.
\end{itemize}

\begin{example}[A casual exchange between two friends]
\leavevmode
\begin{itemize}
  \item[A:] \ja{きのう、なにした？} \emph{Kinō, nani shita?}
        What did you do yesterday?
  \item[B:] \ja{ともだちとえいがをみた。すごくおもしろかったよ。}
        \emph{Tomodachi to eiga o mita. Sugoku omoshirokatta yo.}
        I saw a movie with a friend. It was really interesting.
  \item[A:] \ja{いいなあ。わたしはぜんぜんでかけなかった。}
        \emph{Ii nā. Watashi wa zenzen dekakenakatta.}
        Lucky you. I didn't go out at all.
  \item[B:] \ja{じゃあ、こんどいっしょにいかない？} \emph{Jā, kondo issho ni ikanai?}
        Then do you want to go together next time?
  \item[A:] \ja{うん、いく！} \emph{Un, iku!} Yeah, I'll go!
\end{itemize}
\end{example}

\begin{aside}
\textbf{When to switch.} Japanese students meeting for the first time often begin in
\ja{です}/\ja{ます} and switch to casual speech once they establish they are the same
age or become friends. Speaking casually to a teacher or older stranger sounds rude;
staying in \ja{です}/\ja{ます} with a close friend sounds distant. Exchange students
are usually forgiven, but it is a good idea to follow the other person's lead. The
\ja{せんぱい}/\ja{こうはい} (senior/junior) relationship in clubs means even a one-year
age difference calls for polite speech from the junior.
\end{aside}

\section{Short forms inside sentences}

\subsection{\ja{〜とおもいます} ``I think that\ldots''}
\textbf{short form + \ja{とおもいます}}
\begin{example}
\leavevmode
\begin{itemize}
  \item \ja{あしたはあめがふるとおもいます。} I think it will rain tomorrow.
  \item \ja{やまださんはこないとおもいます。} I think Yamada-san isn't coming.
  \item \ja{あのひとはせんせいだとおもいます。} I think that person is a teacher.
        (\ja{だ} is required here.)
\end{itemize}
\end{example}

\subsection{\ja{〜といっていました} ``(someone) said that\ldots''}
\ja{せんせいはあしたテストがあるといっていました。} The teacher said there's a test
tomorrow.

\subsection{\ja{〜でしょう}/\ja{〜でしょう？} ``probably; right?''}
\ja{あしたはさむいでしょう。} It will probably be cold tomorrow (weather-forecast style).
\ja{これ、おいしいでしょう？} This is good, right?

\section{Kanji}
\begin{center}
\begin{tabular}{lllll}
\toprule
Kanji & Meaning & Readings & Example & \\
\midrule
\ja{学} & study & \ja{ガク}; \ja{まな(ぶ)} & \ja{大学} & daigaku, university \\
\ja{生} & life, birth & \ja{セイ}; \ja{い(きる)}, \ja{う(まれる)} & \ja{学生} & gakusei, student \\
\ja{先} & ahead, previous & \ja{セン}; \ja{さき} & \ja{先生} & sensei, teacher \\
\ja{友} & friend & \ja{ユウ}; \ja{とも} & \ja{友だち} & tomodachi, friend \\
\ja{何} & what & \ja{カ}; \ja{なに}, \ja{なん} & \ja{何時} & nanji, what time \\
\ja{今} & now & \ja{コン}; \ja{いま} & \ja{今日} & kyō, today (special reading) \\
\ja{時} & time, hour & \ja{ジ}; \ja{とき} & \ja{三時} & sanji, 3 o'clock \\
\ja{分} & minute; divide & \ja{フン}, \ja{ブン}; \ja{わ(かる)} & \ja{分かります} & wakarimasu, understand \\
\bottomrule
\end{tabular}
\end{center}

\begin{remark}
\ja{今日} ``today'' is read \ja{きょう}, a special reading (\emph{jukujikun}) for the
whole word that does not come from the individual kanji. Others we know: \ja{明日}
\ja{あした} ``tomorrow'', \ja{昨日} \ja{きのう} ``yesterday'', \ja{一人} \ja{ひとり},
\ja{二人} \ja{ふたり}, \ja{大人} \ja{おとな} ``adult''.
\end{remark}

\section{End-of-term review}

\begin{center}
\begin{tikzpicture}[node distance=0.4cm and 0.5cm,
    hub/.style={draw, circle, fill=blue!10, font=\small\bfseries, align=center, minimum size=2.2cm},
    t/.style={draw, rounded corners, fill=gray!8, font=\scriptsize, align=center, text width=2.6cm, minimum height=0.9cm}]
  \node[hub] (c) at (0,0) {JAPAN\\101R};
  \node[t] (a) at (90:2.9) {kana: \ja{ひらがな}, \ja{カタカナ}, about 45 kanji};
  \node[t] (b) at (30:3.4) {\ja{XはYです}, \ja{の}, \ja{か}, \ja{も}};
  \node[t] (d) at (-30:3.4) {ko-so-a-do; \ja{あります}/\ja{います}; location};
  \node[t] (e) at (-90:2.9) {verbs: \ja{ます}, past, \ja{て}-form, short forms};
  \node[t] (f) at (-150:3.4) {particles: \ja{は が を に へ で と も}};
  \node[t] (g) at (150:3.4) {adjectives: \ja{い}/\ja{な}, \ja{すき}/\ja{きらい}};
  \foreach \n in {a,b,d,e,f,g} \draw (c) -- (\n);
\end{tikzpicture}
\end{center}

\section*{Lesson 12 review}
\begin{itemize}
  \item Short forms: \ja{よむ}/\ja{よまない}/\ja{よんだ}/\ja{よまなかった};
        \ja{い}-adj drop \ja{です}; \ja{な}-adj and nouns use \ja{だ}/\ja{だった}.
  \item \ja{ない}-form: \ja{る}-verbs drop \ja{る}; \ja{う}-verbs \emph{u}$\to$\emph{a}
        (\ja{う}$\to$\ja{わ}); \ja{しない}, \ja{こない}, \ja{ある}$\to$\ja{ない}.
  \item \ja{た}-form = \ja{て}-form with \ja{て}/\ja{で}$\to$\ja{た}/\ja{だ}.
  \item Casual speech: short forms, no \ja{か}, dropped particles, \ja{うん}/\ja{ううん}.
  \item Short form + \ja{とおもいます} / \ja{といっていました}.
  \item Kanji: \ja{学生先友何今時分}.
\end{itemize}
